\documentclass[11pt, a4paper]{article}

\usepackage[T1]{fontenc}
\usepackage[utf8]{inputenc}
\usepackage{lmodern}
\usepackage{microtype}
\usepackage{amsmath, amssymb, amsthm}
\usepackage{mathtools}
\usepackage{geometry}
\usepackage{xcolor}
\definecolor{teal}{RGB}{0,120,120}
\usepackage{titlesec}
\usepackage{titletoc}
\usepackage{booktabs}
\usepackage{array}
\usepackage{listings}
\usepackage[colorlinks=true, linkcolor=teal, citecolor=teal, urlcolor=teal]{hyperref}
\usepackage{cleveref}

\geometry{margin=2.5cm}

\definecolor{codebg}{RGB}{248,248,248}

\titleformat{\section}{\large\bfseries\color{teal}}{{\thesection}}{1em}{}[{\color{teal}\titlerule[0.4pt]}]
\titleformat{\subsection}{\normalsize\bfseries\color{teal}}{{\thesubsection}}{1em}{}
\titleformat{\subsubsection}{\normalsize\itshape}{\thesubsubsection}{1em}{}

\lstset{
  backgroundcolor=\color{codebg},
  basicstyle=\ttfamily\small,
  breaklines=true,
  frame=single,
  rulecolor=\color{teal!40},
  keywordstyle=\color{teal}\bfseries,
  commentstyle=\color{gray},
  language=Python
}

\newtheorem{definition}{Definition}[section]
\newtheorem{theorem}{Theorem}[section]
\newtheorem{proposition}{Proposition}[section]
\newtheorem{corollary}{Corollary}[section]
\newtheorem{remark}{Remark}[section]

\newcommand{\R}{\mathbb{R}}
\newcommand{\Z}{\mathbb{Z}}
\newcommand{\N}{\mathbb{N}}
\newcommand{\Sym}{\mathrm{Sym}}
\newcommand{\FD}{\mathrm{FD}}
\newcommand{\Orb}{\mathrm{Orb}}
\newcommand{\Stab}{\mathrm{Stab}}
\newcommand{\card}[1]{\lvert #1 \rvert}
\newcommand{\norm}[1]{\lVert #1 \rVert}
\newcommand{\smin}{\mathbin{\tilde{\wedge}}}
\newcommand{\smax}{\mathbin{\tilde{\vee}}}
\newcommand{\Td}{T_d}
\newcommand{\Oh}{O_h}
\newcommand{\Ih}{I_h}

\title{%
  \textcolor{teal}{\rule{\linewidth}{0.8pt}}\\[0.4em]
  \textbf{G-Invariant SDF Grammars for Generative Metal Form}\\
  {\large Mathematical Foundation, Genetic Evolution, and Fabrication Constraints}\\[0.3em]
  \textcolor{teal}{\rule{\linewidth}{0.8pt}}
}
\author{SierpSphere Project — \textit{next generation}}
\date{}

\begin{document}

\maketitle
\tableofcontents
\vspace{2em}

% ─────────────────────────────────────────────────────────────────────────────
\section{Finite Symmetry Groups of the Platonic Solids}
\label{sec:groups}
% ─────────────────────────────────────────────────────────────────────────────

\subsection{The Five Platonic Solids and Their Symmetry Groups}

The five Platonic solids are the unique convex polyhedra whose faces, edges
and vertices are all congruent regular polygons meeting at equal angles. Group
theory classifies their symmetries as finite subgroups of the orthogonal group
$O(3)$.

\begin{definition}[Symmetry group of a solid]
The \emph{symmetry group} $G(P)$ of a polyhedron $P$ is the group of all
isometries of $\R^3$ (rotations and reflections) that map $P$ to itself.
\end{definition}

The five Platonic solids produce exactly three distinct symmetry groups:

\begin{table}[h]
\centering
\begin{tabular}{lllcc}
\toprule
Solid & Faces & Symmetry group & Order & Dual \\
\midrule
Tetrahedron  & 4 triangles   & $\Td$  & 24  & self-dual \\
Cube         & 6 squares     & $\Oh$  & 48  & octahedron \\
Octahedron   & 8 triangles   & $\Oh$  & 48  & cube \\
Dodecahedron & 12 pentagons  & $\Ih$  & 120 & icosahedron \\
Icosahedron  & 20 triangles  & $\Ih$  & 120 & dodecahedron \\
\bottomrule
\end{tabular}
\caption{Platonic solids, symmetry groups, and dual pairs.}
\label{tab:platonic}
\end{table}

The \emph{dual} of a Platonic solid is obtained by replacing vertices with
face centres, and vice versa. Duals share the same symmetry group. Therefore
five solids yield exactly three groups: $\Td$, $\Oh$, $\Ih$.

\subsection{Structure of Each Group}

\paragraph{Tetrahedral group $\Td$ (order 24).}
Contains 12 proper rotations (the chiral tetrahedral group $T$) and 12
improper isometries (reflections and rotoreflections). Generators: a $C_3$
rotation about a vertex axis and a $C_2$ rotation about an edge midpoint axis.
The four $C_3$ axes correspond to the four vertices of the tetrahedron.

\paragraph{Octahedral group $\Oh$ (order 48).}
Contains 24 proper rotations (chiral octahedral group $O$) and 24 improper
isometries. Generators: a $C_4$ rotation about a face-centre axis and a $C_3$
rotation about a vertex axis. The three $C_4$ axes align with the coordinate
axes $x, y, z$.

\paragraph{Icosahedral group $\Ih$ (order 120).}
Contains 60 proper rotations (chiral icosahedral group $I$) and 60 improper
isometries. The group $I$ is isomorphic to $A_5$, the alternating group on
five elements—the unique simple group of order 60. The six $C_5$ axes connect
opposite vertex pairs. The golden ratio $\varphi = (1+\sqrt{5})/2$ appears
throughout the coordinate geometry.

\subsection{Fundamental Domains}
\label{sec:fd}

\begin{definition}[Fundamental domain]
A \emph{fundamental domain} $\FD(G)$ for the action of $G$ on $S^2$ is a
closed region $F \subset S^2$ such that:
\begin{enumerate}
  \item $S^2 = \bigcup_{g \in G} g \cdot F$ \quad (tiles the sphere),
  \item $\mathrm{int}(F) \cap \mathrm{int}(g \cdot F) = \emptyset$ for all $g \neq e$ \quad (non-overlapping interiors).
\end{enumerate}
\end{definition}

For finite subgroups of $O(3)$, the fundamental domain is a \emph{spherical
triangle} — a region bounded by three great-circle arcs. Its area is
$4\pi / \card{G}$:

\begin{table}[h]
\centering
\begin{tabular}{llcl}
\toprule
Group & Order & FD area & Corner vectors (unit, approximate) \\
\midrule
$\Td$ & 24  & $\pi/6$  & $(1,1,1)/\sqrt{3}$,\; $(1,1,0)/\sqrt{2}$,\; $(1,0,0)$ \\
$\Oh$ & 48  & $\pi/12$ & $(1,1,1)/\sqrt{3}$,\; $(1,1,0)/\sqrt{2}$,\; $(1,0,0)$ \\
$\Ih$ & 120 & $\pi/30$ & $(1,1,1)/\sqrt{3}$,\; $(\varphi,1,0)/\sqrt{\varphi^2+1}$,\; $(1,0,0)$ \\
\bottomrule
\end{tabular}
\caption{Fundamental domains on $S^2$. The $\Td$ and $\Oh$ share corner directions
but differ in the great-circle arcs bounding them.}
\label{tab:fd}
\end{table}

A point in $\FD(G)$ is parameterised by barycentric coordinates $(u,v)$ with
$u,v \geq 0$, $u+v \leq 1$:
\begin{equation}
  \varphi(u,v) = \frac{(1-u-v)\,\mathbf{c}_0 + u\,\mathbf{c}_1 + v\,\mathbf{c}_2}
                      {\norm{(1-u-v)\,\mathbf{c}_0 + u\,\mathbf{c}_1 + v\,\mathbf{c}_2}}
  \label{eq:fdpoint}
\end{equation}
where $\mathbf{c}_0, \mathbf{c}_1, \mathbf{c}_2$ are the three corner unit vectors of
$\FD(G)$ from \cref{tab:fd}.

% ─────────────────────────────────────────────────────────────────────────────
\section{G-Invariant Signed Distance Functions}
\label{sec:sdf}
% ─────────────────────────────────────────────────────────────────────────────

\subsection{Signed Distance Functions}

\begin{definition}[SDF]
A \emph{signed distance function} is a map $f: \R^3 \to \R$ such that
$f(\mathbf{x}) = 0$ defines an implicit surface $\Sigma$, $f(\mathbf{x}) < 0$
for interior points, and $f(\mathbf{x}) > 0$ for exterior points.
Optionally, $\norm{\nabla f} = 1$ (Eikonal condition) makes $|f|$ the
Euclidean distance to $\Sigma$.
\end{definition}

\begin{definition}[G-invariant SDF]
An SDF $f$ is \emph{G-invariant} (or \emph{G-equivariant under the trivial
representation}) if
\begin{equation}
  f(g \cdot \mathbf{x}) = f(\mathbf{x}) \quad \forall\, g \in G,\; \mathbf{x} \in \R^3.
  \label{eq:invariant}
\end{equation}
The zero set of a G-invariant SDF is a G-symmetric surface.
\end{definition}

\subsection{Canonical G-Primitive SDFs}
\label{sec:primitives}

For each group $G$, the canonical primitive is the SDF of the associated
Platonic solid. These are the \emph{unique} SDFs that are G-invariant and
have exactly G as their symmetry group (no more, no less).

\paragraph{Tetrahedron SDF ($\Td$-invariant).}
Following Quilez~\cite{quilez_sdf}, a regular tetrahedron of inradius $r$:
\begin{equation}
  f_{\mathrm{tet}}(\mathbf{x}) = \frac{1}{\sqrt{3}}\Bigl(
    \max\bigl(-x{-}y{-}z,\; x{+}y{-}z,\; {-}x{+}y{+}z,\; x{-}y{+}z\bigr) - r
  \Bigr)
  \label{eq:tet}
\end{equation}

\paragraph{Cube SDF ($\Oh$-invariant).}
\begin{equation}
  f_{\mathrm{box}}(\mathbf{x}) = \norm{\max(\lvert\mathbf{x}\rvert - \mathbf{h},\, \mathbf{0})}_2
  + \min\bigl(\max(|\mathbf{x}| - \mathbf{h}),\, 0\bigr)
  \label{eq:box}
\end{equation}
where $\mathbf{h} = (r, r, r)$ for a cube of half-extent $r$.

\paragraph{Icosahedron SDF ($\Ih$-invariant).}
Let $\varphi = (1+\sqrt{5})/2$. The icosahedron is the intersection of five
``cube-like'' half-space families. The face normals of its dual (the
dodecahedron) are:
\begin{equation}
  \mathbf{n}_1 = \tfrac{1}{\sqrt{\varphi^2+1}}(\varphi, 1, 0), \quad
  \mathbf{n}_2 = \tfrac{1}{\sqrt{\varphi^2+1}}(0, \varphi, 1), \quad
  \mathbf{n}_3 = \tfrac{1}{\sqrt{\varphi^2+1}}(1, 0, \varphi)
\end{equation}
and all sign permutations under $\Ih$. The SDF is:
\begin{equation}
  f_{\mathrm{ico}}(\mathbf{x}) = \max_{i}\bigl(\lvert \mathbf{n}_i \cdot \mathbf{x} \rvert\bigr)
  \cdot \frac{\sqrt{\varphi^2+1}}{\varphi} - r
  \label{eq:ico}
\end{equation}

\subsection{The Orbit Symmetrisation Operator}
\label{sec:symmetrize}

To add a child primitive at position $\boldsymbol{\phi} \in \FD(G)$ while
preserving G-invariance, we define the \emph{orbit symmetrisation}:

\begin{definition}[Symmetrize$_G$]
\begin{equation}
  \mathrm{Sym}_G(P, \boldsymbol{\phi})(\mathbf{x})
    = \min_{g \in G}\, P(g^{-1} \cdot \mathbf{x} - \boldsymbol{\phi})
  \label{eq:symg}
\end{equation}
where $P: \R^3 \to \R$ is any SDF primitive.
\end{definition}

\begin{proposition}[G-invariance of Sym$_G$]
$\mathrm{Sym}_G(P, \boldsymbol{\phi})$ is G-invariant for any $P$ and any
$\boldsymbol{\phi} \in \R^3$.
\end{proposition}
\begin{proof}
For any $h \in G$:
$\mathrm{Sym}_G(P, \boldsymbol{\phi})(h \cdot \mathbf{x})
= \min_{g \in G} P(g^{-1} h \mathbf{x} - \boldsymbol{\phi})
= \min_{g' \in G} P(g'^{-1} \mathbf{x} - \boldsymbol{\phi})
= \mathrm{Sym}_G(P, \boldsymbol{\phi})(\mathbf{x})$,
where we used the substitution $g' = h^{-1}g$ and the fact that left
multiplication is a bijection on $G$. \qed
\end{proof}

The orbit of $\boldsymbol{\phi}$ under $G$ has size $\card{G}/\card{\Stab_G(\boldsymbol{\phi})}$
by the orbit-stabiliser theorem. For a generic point in the interior of
$\FD(G)$, the stabiliser is trivial, so the orbit has size $\card{G}$:
24, 48, or 120 child copies respectively.

\subsection{Smooth Boolean Algebra}

The smooth Boolean operations of Inigo Quilez~\cite{quilez_smooth} are:
\begin{align}
  f \smin_k g &= \mathrm{smin}(f, g, k)
    \coloneqq \mathrm{lerp}(g, f, h) - k\,h\,(1-h),
    \quad h = \mathrm{clamp}\!\left(\tfrac{1}{2} + \tfrac{g-f}{2k}, 0, 1\right)
    \label{eq:smin}\\
  f \smax_k g &= \mathrm{smax}(f, g, k)
    \coloneqq \mathrm{lerp}(g, f, h) + k\,h\,(1-h),
    \quad h = \mathrm{clamp}\!\left(\tfrac{1}{2} - \tfrac{g-f}{2k}, 0, 1\right)
    \label{eq:smax}
\end{align}

\begin{proposition}[Closure under G-invariance]
If $f$ and $g$ are G-invariant SDFs, then $f \smin_k g$ and $f \smax_k g$
are also G-invariant for any $k \geq 0$.
\end{proposition}
\begin{proof}
Both operations are pointwise functions of $f(\mathbf{x})$ and $g(\mathbf{x})$.
Since G-invariance is preserved by pointwise composition, the result follows
immediately. \qed
\end{proof}

% ─────────────────────────────────────────────────────────────────────────────
\section{The G-Grammar}
\label{sec:grammar}
% ─────────────────────────────────────────────────────────────────────────────

\subsection{Formal Definition}

\begin{definition}[G-Grammar]
A \emph{G-grammar} is a tuple
\[
  \Gamma = \bigl(G,\; f_{\mathrm{seed}},\; [(b_1, P_1, \phi_1, k_1), \ldots, (b_n, P_n, \phi_n, k_n)]\bigr)
\]
where:
\begin{itemize}
  \item $G \in \{\Td, \Oh, \Ih\}$ is the symmetry group,
  \item $f_{\mathrm{seed}}$ is the canonical G-primitive SDF (\cref{sec:primitives}),
  \item $b_i \in \{\mathrm{union}, \mathrm{subtract}, \mathrm{intersect}\}$ is a Boolean operation,
  \item $P_i$ is any primitive SDF (need not be G-invariant),
  \item $\phi_i = \mathrm{fd\_point}(G, u_i, v_i)$ with $(u_i, v_i) \in [0,1]^2$, $u_i+v_i\leq 1$,
  \item $k_i \geq 0$ is the smooth-blending radius.
\end{itemize}
\end{definition}

\subsection{Evaluation}

The SDF of $\Gamma$ is defined inductively:
\begin{align}
  f_0 &= f_{\mathrm{seed}} \\
  f_i &= f_{i-1} \;\Box_{b_i, k_i}\; \mathrm{Sym}_G(P_i, s_i \cdot \phi_i)
  \label{eq:eval}
\end{align}
where $s_i > 0$ is a scale factor placing the orbit at radius $s_i$ from the
origin, and $\Box_{b,k}$ denotes the appropriate smooth Boolean from
\cref{eq:smin}--\cref{eq:smax}.

\begin{corollary}[Closure]
$f_n$ is G-invariant for any choice of $(b_i, P_i, \phi_i, k_i, s_i)$.
\end{corollary}
This follows by induction from \cref{eq:invariant}, \cref{eq:symg}, and the
closure proposition above.

\subsection{Bijection: Seed \texorpdfstring{$\leftrightarrow$}{<->} Group}

The assignment $G(\mathrm{seed})$ is a bijection:
\begin{equation}
  \mathrm{tetrahedron} \leftrightarrow \Td, \quad
  \mathrm{cube} \leftrightarrow \Oh, \quad
  \mathrm{icosahedron} \leftrightarrow \Ih.
\end{equation}
This means the symmetry group is \emph{not} a free parameter — it is
determined by the seed choice. This eliminates a class of incoherent grammars
(e.g.\ an octahedral seed declared with icosahedral symmetry).

% ─────────────────────────────────────────────────────────────────────────────
\section{Genetic Algorithm on Grammar Space}
\label{sec:ga}
% ─────────────────────────────────────────────────────────────────────────────

\subsection{Search Space}

The space of G-grammars with at most $n$ iterations is:
\begin{equation}
  \mathcal{G}_n = \{(\Td, \Oh, \Ih)\} \times \prod_{i=1}^{n}
    \bigl(\{b\} \times \{P\} \times [0,1]^2 \times \R_{>0} \times \R_{\geq 0}\bigr)^*
\end{equation}
with the length of the product chosen from $\{1, \ldots, n\}$.

The subspace $\mathcal{G}_n$ is a union of manifolds; the continuous parameters
$(u_i, v_i, s_i, k_i)$ form smooth submanifolds, while the discrete
parameters $(b_i, P_i, \mathrm{seed})$ index disconnected components.

\subsection{Representation and Operators}

\paragraph{Representation.}
Each individual is a JSON-serialisable dict (the G-grammar as defined above),
making the population trivially persistent and readable.

\paragraph{Initialisation.}
$\mathrm{diverse\_population}(N)$ seeds $N/3$ individuals from each of the
three seed types, with $n_{\mathrm{steps}} \sim \mathrm{Uniform}(1,4)$,
and all continuous parameters sampled uniformly in their domains.

\paragraph{Mutation.}
With rate $p_m = 0.55$, each gene mutates independently:
\begin{itemize}
  \item $(u_i, v_i) \leftarrow \mathrm{clamp}((u_i + \mathcal{N}(0, 0.15),\; v_i + \mathcal{N}(0, 0.15)),\; [0,1]^2)$, then renormalise if $u_i + v_i > 1$.
  \item $b_i$ resampled from $\{\mathrm{union, subtract, intersect}\}$ uniformly.
  \item $P_i$ resampled from $\{\mathrm{tetrahedron, cube, icosahedron, sphere}\}$.
  \item $s_i, k_i$ perturbed log-normally.
  \item Seed type mutated with probability $0.15 \cdot p_m$ (changes $G$, all $\phi_i$ resampled in new $\FD$).
\end{itemize}

\paragraph{Crossover (relaxed).}
Given parents $\Gamma_A, \Gamma_B$:
\begin{enumerate}
  \item Child $C_A$ inherits seed from $\Gamma_A$ (determines $G_A$); child $C_B$ from $\Gamma_B$.
  \item Operations are drawn freely from either parent's operation list (uniform random index selection).
  \item If an operation borrowed from $\Gamma_B$ carries $(u_i, v_i)$ and $G_A \neq G_B$, the coordinates are resampled uniformly in $\FD(G_A)$ (the FD topology differs; the old coordinates may not map coherently).
\end{enumerate}
This is \emph{group-coherent crossover}: every child has a valid G-grammar
with respect to its inherited group, while benefiting from structural
variation borrowed across groups.

\paragraph{Selection.}
Tournament selection with $k=2$. Low $k$ preserves diversity; type-aware
elitism guarantees no seed type goes extinct:
\begin{itemize}
  \item Top-2 individuals (global elitism),
  \item Best viable individual of each seed type (type elitism).
\end{itemize}

\subsection{Convergence and Diversity}

The combination of weak selection ($k=2$), high mutation ($p_m = 0.55$), type
elitism, and diverse initialisation maintains exploration across the three
disconnected subspaces $\mathcal{G}_{\Td}, \mathcal{G}_{\Oh}, \mathcal{G}_{\Ih}$.

Within each subspace, the continuous parameters $(u,v,s,k)$ form a smooth
manifold on which the fitness landscape is differentiable almost everywhere
(smooth Booleans are $C^\infty$ except at $k=0$). This makes the landscape
amenable to both gradient-free (GA) and potentially gradient-based (CMA-ES)
optimisation.

% ─────────────────────────────────────────────────────────────────────────────
\section{Fitness Landscape for Metal Additive Manufacturing}
\label{sec:fitness}
% ─────────────────────────────────────────────────────────────────────────────

\subsection{Hard Constraints (Manufacturing Gates)}

Three constraints are evaluated as hard gates — violation immediately returns
fitness $= 0$, regardless of aesthetic metrics:

\begin{enumerate}
  \item \textbf{Connectivity}: the mesh must be a single connected component (no floating islands, which cannot be printed as a single part).
  \item \textbf{Watertightness}: the mesh must have no boundary edges (open surfaces trap powder in DMLS).
  \item \textbf{Minimum wall thickness}: $t_{\min} \geq 0.3\,\mathrm{mm}$ at target scale $80\,\mathrm{mm}$. This is the resolution limit of standard DMLS systems with $30\,\mu\mathrm{m}$ laser spot.
\end{enumerate}

\subsection{Fitness Metrics}

All soft metrics $s_i \in [0,1]$. The fitness is:
\begin{equation}
  F(\Gamma) = \sum_{i} w_i \, s_i(\Gamma), \quad \sum_i w_i = 1.
  \label{eq:fitness}
\end{equation}

\begin{table}[h]
\centering
\begin{tabular}{llp{7cm}}
\toprule
Metric & Weight & Description \\
\midrule
fractal\_dimension     & 0.16 & Box-counting dimension of mesh vertices; Gaussian reward centred at 2.5 \\
curvature\_variance    & 0.12 & Std of discrete mean curvature; rewards complex, varied surface \\
normalised\_S/V        & 0.08 & Surface-to-volume ratio normalised by convex hull; rewards surface complexity \\
min\_wall\_thickness   & 0.15 & Score derived from ray-cast minimum thickness \\
min\_feature\_size     & 0.08 & Proxy: same as wall thickness \\
fill\_ratio            & 0.05 & $1 - V_{\mathrm{mesh}}/V_{\mathrm{hull}}$; rewards hollowness/carving \\
genus                  & 0.05 & Topological holes; $g = 1 - \chi/2$; rewards handles \\
aspect\_ratio          & 0.03 & Penalises elongated forms ($> 10:1$) \\
drain\_openings        & 0.05 & Guaranteed by watertight gate; always 1.0 \\
enclosed\_voids        & 0.03 & Guaranteed by gate; always 1.0 \\
no\_islands            & 0.03 & Guaranteed by gate; always 1.0 \\
thermal\_mass\_variance& 0.05 & Octant face-area variance; rewards uniform mass distribution \\
support\_volume\_ratio & 0.05 & Overhang fraction; penalises print supports \\
silhouette\_complexity & 0.05 & Projected area variance across 6 orthogonal views \\
primitive\_diversity   & 0.02 & Distinct primitive types in the grammar \\
\midrule
\textbf{Total}         & \textbf{1.00} & \\
\bottomrule
\end{tabular}
\caption{Fitness metric weights. \texttt{symmetry\_preservation} is absent:
G-invariance is structural (see \cref{sec:grammar}), not measured.}
\label{tab:weights}
\end{table}

\subsection{The Fill Ratio Metric}

\begin{definition}[Fill ratio]
\[
  s_{\mathrm{fill}}(\Gamma) = 1 - \frac{V_{\mathrm{mesh}}}{V_{\mathrm{hull}} + \varepsilon}
\]
where $V_{\mathrm{mesh}}$ is the mesh volume and $V_{\mathrm{hull}}$ is its
convex hull volume.
\end{definition}

A solid blob scores 0 (volume $\approx$ hull); a maximally hollowed lattice
scores approaching 1. This reward is consistent with the manufacturing
objective of minimising material cost and residual thermal stress in DMLS,
while the hard gate on minimum wall thickness prevents degenerate ultra-thin shells.

% ─────────────────────────────────────────────────────────────────────────────
\section{Computational Architecture}
\label{sec:compute}
% ─────────────────────────────────────────────────────────────────────────────

\subsection{GPU-Accelerated SDF Evaluation}

The SDF grid is evaluated on Apple Silicon Metal via PyTorch MPS
\cite{apple_mps}. For a resolution of $N^3$ voxels:

\begin{enumerate}
  \item Build the coordinate tensor $\mathbf{X} \in \R^{N^3 \times 3}$ on the MPS device.
  \item Evaluate $f_{\mathrm{seed}}(\mathbf{X})$ as a single batched operation.
  \item For each operation $i$: compute $\mathrm{Sym}_G(P_i, \phi_i)(\mathbf{X})$ by stacking
        $\card{G}$ primitive evaluations and taking the channel-wise minimum.
  \item Apply smooth Boolean to update $f$.
\end{enumerate}

The orbit symmetrisation for $\Ih$ evaluates $120 \times N^3$ primitive SDFs.
At $N=28$: $120 \times 21{,}952 = 2{,}634{,}240$ evaluations per step — fully
parallelised in a single MPS kernel dispatch. Memory cost: $\approx 10\,\mathrm{MB}$
per individual at float32, negligible on 16+ GB unified memory.

\subsection{Parallel Population Evaluation}

Population evaluation uses Python \texttt{multiprocessing} with the
\texttt{spawn} start method (required on macOS to avoid fork-MPS conflicts).
With $C$ CPU cores and population size $P$, $\min(P, C)$ workers are spawned,
each initialising its own MPS context. Wall time per epoch scales as
$\lceil P / C \rceil \times t_{\mathrm{individual}}$.

% ─────────────────────────────────────────────────────────────────────────────
\section{Metal Additive Manufacturing Constraints}
\label{sec:print}
% ─────────────────────────────────────────────────────────────────────────────

\subsection{Direct Metal Laser Sintering (DMLS) and SLM}

DMLS and Selective Laser Melting (SLM) build parts layer by layer by fusing
metal powder with a laser. Relevant constraints for generative design:

\begin{description}
  \item[Layer thickness] Typically $20$--$60\,\mu\mathrm{m}$. Surfaces
    angled below $\approx 45^\circ$ from horizontal require support structures.
  \item[Minimum feature] Limited by laser spot diameter ($\approx 30$--$100\,\mu\mathrm{m}$) and thermal diffusion. We gate on $0.3\,\mathrm{mm}$.
  \item[Enclosed cavities] Trapped powder cannot be removed post-print; the watertightness gate prevents this structurally.
  \item[Thermal gradients] Uneven mass distribution causes residual stress and warping; the thermal mass variance metric rewards octant uniformity.
  \item[Support removal] G-symmetric overhangs distribute support attachment points evenly, simplifying post-processing.
\end{description}

\subsection{Why G-Symmetric Designs Excel at DMLS}

A G-invariant shape distributes all features symmetrically under the group
action. For $\Ih$ (order 120), this means:
\begin{itemize}
  \item Thermal load during sintering is distributed equally across 120
    equivalent facets — minimal warping.
  \item Support structures (if needed) attach at 120 equivalent points,
    distributing stress.
  \item Post-print measurements at any one canonical point characterise all
    equivalent points — simplified QC.
\end{itemize}
These are \emph{consequences} of group-invariance, not additional engineering
constraints. The grammar is thus not merely an aesthetic choice but a
manufacturing-aware design principle.

% ─────────────────────────────────────────────────────────────────────────────
\begin{thebibliography}{99}

\bibitem{coxeter}
H.~S.~M.\ Coxeter,
\textit{Regular Polytopes}, 3rd ed.
Dover Publications, 1973.

\bibitem{quilez_sdf}
I.~Quilez,
``Modelling with Distance Functions'',
\textit{iquilezles.org}, 2008--2024.
\url{https://iquilezles.org/articles/distfunctions/}

\bibitem{quilez_smooth}
I.~Quilez,
``Smooth Minimum'',
\textit{iquilezles.org}, 2013.
\url{https://iquilezles.org/articles/smin/}

\bibitem{lorensen}
W.~E.\ Lorensen and H.~E.\ Cline,
``Marching Cubes: A High Resolution 3D Surface Construction Algorithm'',
\textit{ACM SIGGRAPH Computer Graphics}, 21(4):163--169, 1987.

\bibitem{holland}
J.~H.\ Holland,
\textit{Adaptation in Natural and Artificial Systems}.
University of Michigan Press, 1975.

\bibitem{deb}
K.~Deb, A.~Pratap, S.~Agarwal, and T.~Meyarivan,
``A Fast and Elitist Multiobjective Genetic Algorithm: NSGA-II'',
\textit{IEEE Transactions on Evolutionary Computation}, 6(2):182--197, 2002.

\bibitem{pytorch}
A.~Paszke et~al.,
``PyTorch: An Imperative Style, High-Performance Deep Learning Library'',
\textit{NeurIPS}, 2019.
\url{https://pytorch.org}

\bibitem{apple_mps}
Apple Inc.,
``Metal Performance Shaders'', 2024.
\url{https://developer.apple.com/metal/}

\bibitem{sethian}
J.~A.\ Sethian,
\textit{Level Set Methods and Fast Marching Methods}.
Cambridge University Press, 1999.

\bibitem{wang_dmls}
D.\ Wang et~al.,
``Research on the Fabricating Quality Optimisation of the Overhanging
Surface in Metal Laser Melting'',
\textit{Applied Surface Science}, 285:534--543, 2013.

\bibitem{deepsdf}
J.~J.\ Park, P.~Florence, J.~Straub, R.~Newcombe, and S.~Lovegrove,
``DeepSDF: Learning Continuous Signed Distance Functions for Shape
Representation'',
\textit{CVPR}, 2019.

\bibitem{a5simple}
M.~A.\ Armstrong,
\textit{Groups and Symmetry}.
Springer, 1988.

\end{thebibliography}

\end{document}
